% GENERATED FILE. Edit canonical lessons, then run npm run generate:books.
% canonical-source-hash: fnv1a64:d71238d5474bb4e5
% canonical-lessons: PT-C450-mercearia, PT-C450-charcutaria, PT-C450-frutaria, PT-C450-garrafeira, PT-C450-fatiar

\chapter{Comida na loja}
\label{ch:pt-comida-na-loja}
\hlchaptermodality{450}

\begin{chapteropening}
\textbf{By the end of this chapter:} \emph{I can shop for food: the grocer's, the deli, the fruit shop, the wine shop and slicing.}

Complete the last lesson of chapter 450: I can shop for food: the grocer's, the deli, the fruit shop, the wine shop and slicing.
\end{chapteropening}

\section[a mercearia]{a mercearia — the grocer's, the corner shop}

\label{lesson:PT-C450-mercearia}



\begin{quote}
Before the new one: say the Portuguese for the herbalist's shop, then the Portuguese for the jeweller's shop.
\end{quote}

\subsection*{You'll want to know: a mercearia}
\textbf{a mercearia} — \textquotedblleft{}the grocer's, the corner shop\textquotedblright{}.

\begin{quote}
A mercearia do bairro abre ao domingo.
\end{quote}

\begin{quote}
\emph{The local grocer's opens on Sundays.}
\end{quote}

In Brazil: \emph{a mercearia} or \emph{o armazém}.

\begin{cousinweb}
EtymDB links \emph{mercearia} with Italian \emph{merceria}; the family goes back to Latin \emph{merx, mercis} 'goods' (Lewis \& Short).
\end{cousinweb}

\begin{grammarlens}[title={What you've built}]
One more word for telling, explaining and describing at B1.
\end{grammarlens}

\subsection*{Your turn}
\begin{itemize}
\raggedright
  \item \emph{Say it:} \emph{a mercearia}
  \item \emph{Say it:} \emph{a mercearia}, once more
  \item \emph{Say it:} a sentence of your own with \emph{a mercearia}
  \item \emph{Recall:} say the Portuguese for the herbalist's shop, then the Portuguese for the jeweller's shop, then say \emph{a mercearia} again
\end{itemize}

\begin{tcolorbox}[breakable,colback=teal!4,colframe=teal!35!black,title={Before you move on}]
What does a mercearia mean? (\textquotedblleft{}The grocer's, the corner shop\textquotedblright{}.) Say it once more.
\end{tcolorbox}
\section[a charcutaria]{a charcutaria — the deli (cured meats and cheese)}

\label{lesson:PT-C450-charcutaria}



\begin{quote}
Before the new one: say the Portuguese for the jeweller's shop, then the Portuguese for the grocer's.
\end{quote}

\subsection*{You'll want to know: a charcutaria}
\textbf{a charcutaria} — \textquotedblleft{}the deli (cured meats and cheese)\textquotedblright{}.

\begin{quote}
Na charcutaria compro presunto e chouriço.
\end{quote}

\begin{quote}
\emph{At the deli I buy cured ham and chouriço.}
\end{quote}

Also the deli counter in a supermarket.

\begin{cousinweb}
\emph{charcutaria} was borrowed from French \emph{charcuterie}, from \emph{chair cuite} 'cooked flesh'; \emph{chair} is inherited from Latin \emph{carō} 'flesh' (Lewis \& Short).
\end{cousinweb}

\begin{grammarlens}[title={What you've built}]
One more word for telling, explaining and describing at B1.
\end{grammarlens}

\subsection*{Your turn}
\begin{itemize}
\raggedright
  \item \emph{Say it:} \emph{a charcutaria}
  \item \emph{Say it:} \emph{a charcutaria}, once more
  \item \emph{Say it:} a sentence of your own with \emph{a charcutaria}
  \item \emph{Recall:} say the Portuguese for the jeweller's shop, then the Portuguese for the grocer's, then say \emph{a charcutaria} again
\end{itemize}

\begin{tcolorbox}[breakable,colback=teal!4,colframe=teal!35!black,title={Before you move on}]
What does a charcutaria mean? (\textquotedblleft{}The deli (cured meats and cheese)\textquotedblright{}.) Say it once more.
\end{tcolorbox}
\section[a frutaria]{a frutaria — the fruit shop, the greengrocer}

\label{lesson:PT-C450-frutaria}



\begin{quote}
Before the new one: say the Portuguese for the grocer's, then the Portuguese for the deli.
\end{quote}

\subsection*{You'll want to know: a frutaria}
\textbf{a frutaria} — \textquotedblleft{}the fruit shop, the greengrocer\textquotedblright{}.

\begin{quote}
Na frutaria, as laranjas estão baratas.
\end{quote}

\begin{quote}
\emph{At the fruit shop, oranges are cheap.}
\end{quote}

Built on \emph{a fruta}.

\begin{cousinweb}
\emph{frutaria} is built on \emph{fruta} 'fruit', which goes back to Latin \emph{frūctus} 'enjoyment, fruit' (Lewis \& Short), through a Late Latin plural form \emph{frūcta}.
\end{cousinweb}

\begin{grammarlens}[title={What you've built}]
One more word for telling, explaining and describing at B1.
\end{grammarlens}

\subsection*{Your turn}
\begin{itemize}
\raggedright
  \item \emph{Say it:} \emph{a frutaria}
  \item \emph{Say it:} \emph{a frutaria}, once more
  \item \emph{Say it:} a sentence of your own with \emph{a frutaria}
  \item \emph{Recall:} say the Portuguese for the grocer's, then the Portuguese for the deli, then say \emph{a frutaria} again
\end{itemize}

\begin{tcolorbox}[breakable,colback=teal!4,colframe=teal!35!black,title={Before you move on}]
What does a frutaria mean? (\textquotedblleft{}The fruit shop, the greengrocer\textquotedblright{}.) Say it once more.
\end{tcolorbox}
\section[a garrafeira]{a garrafeira — the wine shop; the wine cellar}

\label{lesson:PT-C450-garrafeira}



\begin{quote}
Before the new one: say the Portuguese for the deli, then the Portuguese for the fruit shop.
\end{quote}

\subsection*{You'll want to know: a garrafeira}
\textbf{a garrafeira} — \textquotedblleft{}the wine shop; the wine cellar\textquotedblright{}.

\begin{quote}
Comprámos um vinho velho na garrafeira.
\end{quote}

\begin{quote}
\emph{We bought an old wine at the wine shop.}
\end{quote}

\emph{Vinho garrafeira} — a long-aged wine.

\begin{cousinweb}
\emph{garrafeira} is built on \emph{garrafa} 'bottle' (EtymDB) with \emph{-eira}, a place for them; the origin of \emph{garrafa} is not settled.
\end{cousinweb}

\begin{grammarlens}[title={What you've built}]
One more word for telling, explaining and describing at B1.
\end{grammarlens}

\subsection*{Your turn}
\begin{itemize}
\raggedright
  \item \emph{Say it:} \emph{a garrafeira}
  \item \emph{Say it:} \emph{a garrafeira}, once more
  \item \emph{Say it:} a sentence of your own with \emph{a garrafeira}
  \item \emph{Recall:} say the Portuguese for the deli, then the Portuguese for the fruit shop, then say \emph{a garrafeira} again
\end{itemize}

\begin{tcolorbox}[breakable,colback=teal!4,colframe=teal!35!black,title={Before you move on}]
What does a garrafeira mean? (\textquotedblleft{}The wine shop; the wine cellar\textquotedblright{}.) Say it once more.
\end{tcolorbox}
\section[fatiar]{fatiar — to slice}

\label{lesson:PT-C450-fatiar}



\begin{quote}
Before the new one: say the Portuguese for the fruit shop, then the Portuguese for the wine shop.
\end{quote}

\subsection*{You'll want to know: fatiar}
\textbf{fatiar} — \textquotedblleft{}to slice\textquotedblright{}.

\begin{quote}
O senhor fatia o fiambre muito fino.
\end{quote}

\begin{quote}
\emph{The man slices the ham very thin.}
\end{quote}

Built on \emph{a fatia}, the slice.

\begin{cousinweb}
\emph{fatiar} is built on \emph{fatia} 'slice', which EtymDB traces to Arabic \emph{futāt} 'crumbs, pieces'.
\end{cousinweb}

\begin{grammarlens}[title={What you've built}]
One more word for telling, explaining and describing at B1.
\end{grammarlens}

\subsection*{Your turn}
\begin{itemize}
\raggedright
  \item \emph{Say it:} \emph{fatiar}
  \item \emph{Say it:} \emph{fatiar}, once more
  \item \emph{Say it:} a sentence of your own with \emph{fatiar}
  \item \emph{Recall:} say the Portuguese for the fruit shop, then the Portuguese for the wine shop, then say \emph{fatiar} again
\end{itemize}

\begin{tcolorbox}[breakable,colback=teal!4,colframe=teal!35!black,title={Before you move on}]
What does fatiar mean? (\textquotedblleft{}To slice\textquotedblright{}.) Say it once more.
\end{tcolorbox}
